%% kapa: exam=1 round=28 year=2017 period=2 form=A subject=appraisal_law subject_ko=감정평가관계법규 seq=7 number=7
%% source: Q-Net 감정평가사 기출문제 (cst00309) artlSeq=5199603 file=extracted/1008/5199603/2017년 제28회 감정평가사 1차시험 문제지/2교시 A형.pdf
%% conversion: kordoc 4.24.51 -> tools/gen1.py (auto); review=proofread-sonnet; flags=none
\begin{question}{7}
국토의 계획 및 이용에 관한 법령상 지구단위계획구역에 관한 설명으로 옳지 않은 것은?
\begin{choices}
  \item 지구단위계획구역은 도시ㆍ군관리계획으로 결정한다.
  \item 용도지구로 지정된 지역에 대하여는 지구단위계획구역을 지정할 수 없다.
  \item 「도시 및 주거환경정비법」에 따라 지정된 정비구역의 일부에 대하여 지구단위계획구역을 지정할 수 있다.
  \item 도시지역 외 지구단위계획구역에서는 지구단위계획으로 당해 용도지역 또는 개발진흥지구에 적용되는 건폐율의 150 퍼센트 이내에서 건폐율을 완화하여 적용할 수 있다.
  \item 도시지역 내 지구단위계획구역의 지정목적이 한옥마을을 보존하고자 하는 경우 지구단위계획으로「주차장법」에 의한 주차장 설치기준을 100 퍼센트까지 완화하여 적용할 수 있다.
\end{choices}
\end{question}
\answer{2}
